\documentclass[10pt,twocolumn]{article}
\usepackage[letterpaper,margin=0.75in,columnsep=0.25in]{geometry}
\usepackage{fontspec}
\setmainfont{DejaVu Serif}
\setsansfont{DejaVu Sans}
\setmonofont{DejaVu Sans Mono}
\usepackage{amsmath,amssymb}
\usepackage{graphicx}
\usepackage{xcolor}
\usepackage{listings}
\usepackage{hyperref}
\usepackage{caption}
\usepackage{titlesec}
\usepackage{authblk}
\usepackage{enumitem}

% --- IEEE-like section formatting ---
\titleformat{\section}{\centering\normalfont\small\bfseries}{\Roman{section}.}{0.5em}{\MakeUppercase}
\titleformat{\subsection}{\normalfont\small\bfseries\itshape}{\Alph{subsection}.}{0.5em}{}
\titlespacing*{\section}{0pt}{12pt}{6pt}
\titlespacing*{\subsection}{0pt}{8pt}{4pt}

\renewcommand{\Authfont}{\normalfont\normalsize}
\renewcommand{\Affilfont}{\normalfont\small\itshape}
\setlength{\affilsep}{4pt}

% --- Code listing colors ---
\definecolor{codebg}{RGB}{245,245,245}
\definecolor{codekw}{RGB}{0,0,150}
\definecolor{codecomment}{RGB}{100,100,100}
\definecolor{codestring}{RGB}{160,20,20}

% Example language definition (Go). Replace `keywords` with the
% appropriate set for your topic's language.
\lstdefinelanguage{Go}{
  keywords={package, import, func, type, struct, interface, return, if, else, for, range, var, const, defer, go, chan, select, case, switch, default, nil, true, false, map, break, continue, err, panic, recover},
  keywordstyle=\color{codekw}\bfseries,
  ndkeywords={string, int, bool, error, float64, byte, rune},
  ndkeywordstyle=\color{teal}\bfseries,
  sensitive=true,
  comment=[l]{//},
  morecomment=[s]{/*}{*/},
  commentstyle=\color{codecomment}\itshape,
  stringstyle=\color{codestring},
  morestring=[b]",
  morestring=[b]`,
}

\lstset{
  language=Go,
  backgroundcolor=\color{codebg},
  basicstyle=\fontsize{7}{8}\ttfamily,
  breaklines=true,
  frame=single,
  numbers=left,
  numberstyle=\tiny\color{gray},
  showstringspaces=false,
  tabsize=2,
  xleftmargin=1em,
  captionpos=b,
  aboveskip=6pt,
  belowskip=2pt
}

\pagestyle{plain}

\title{\bfseries PAPER TITLE GOES HERE (any language / Unicode)}
\author{AUTHOR NAME}
\affil{AFFILIATION / DATE LINE}
\date{}

\begin{document}
\twocolumn[
  \begin{@twocolumnfalse}
    \maketitle
    \begin{abstract}
    \small
    \noindent ABSTRACT TEXT in the target language: 150-250 words.

    \vspace{4pt}
    \noindent\textit{Index Terms}---keyword1, keyword2, keyword3.
    \end{abstract}
    \vspace{10pt}
  \end{@twocolumnfalse}
]

\section{Introduction}
Section body text in the target language.

\section{Background}

\subsection{First subsection}
Body text.

\begin{lstlisting}[caption={Caption text.}]
// example code — code/comments conventionally stay in English
// even when the surrounding prose is in another language
func example() {
    return
}
\end{lstlisting}

\section{Main Content}
Core content.

\begin{table}[htbp]
\centering
\small
\caption{Table caption.}
\label{tab:example}
\begin{tabular}{p{1.6cm}p{2.3cm}p{2.2cm}}
\hline
\textbf{Aspect} & \textbf{A} & \textbf{B} \\
\hline
Row & value & value \\
\hline
\end{tabular}
\end{table}

\section{Conclusion}
Closing summary.

\begin{thebibliography}{9}
\small
\bibitem{ref1} Author, ``Title,'' Source, Year.
\end{thebibliography}

\end{document}
